%%
%% fall-lecture02.tex
%% 
%% Made by Alex Nelson <pqnelson@gmail.com>
%% Login   <alex@lisp>
%% 
%% Started on  2026-10-04T12:45:16-0700
%% Last update 2026-10-04T12:45:16-0700
%% 

\lecture{}

\subsection{Continuous maps}

\begin{definition}
Let $(X,\tau_{X})$, $(Y,\tau_{Y})$ be topological spaces.
Let $f\colon X\to Y$ be a function. Let $x_{0}\in X$ and
$y_{0}=f(x_{0})\in Y$. We say that $f$ is \define{Continuous at the Point}
$x_{0}$ if for any neighborhood $V\subset Y$ of $y_{0}\in V$ there
exists a neighborhood $x_{0}\in U\subset X$ such that $f(U)\subset V$.
\end{definition}

\begin{example}
If $X$ and $Y$ are metric spaces, then for every $\varepsilon>0$ there
is a $\delta>0$ such that
\begin{equation}
\forall x\in X\ldotp d(x,x_{0})<\delta\implies \rho(f(x),f(x_{0}))<\varepsilon
\end{equation}
is the usual notion of continuity.
\end{example}

\begin{definition}
We say $f\colon X\to Y$ is \define{Continuous} if $f$ is continuous at
every point of $X$.
\end{definition}

\begin{theorem}
We have $f$ is (globally) continuous if and only if the pre-image of
every open set is open: $V\subset Y$ is open implies $f^{*}V\subset X$
is open.
\end{theorem}

\subsection{Product topology}

\begin{node}
We have a family $(Y_{\lambda},\tau_{\lambda})$ of topological spaces
where $\lambda\in\Lambda$ is an element of an indexing set. We can
consider the \emph{set}
\begin{equation}
Y=\prod_{\lambda\in\Lambda}Y_{\lambda}.
\end{equation}
Elements of $Y$ are generalized sequences
\begin{equation}
\vec{y}=(y_{\lambda})_{\lambda\in\Lambda}
\end{equation}
where $y_{\lambda}\in Y_{\lambda}$. We write $\vec{Y}=\prod_{\lambda\in\Lambda}Y_{\lambda}$. Now, we have described the set, we
want to describe the topology.

We denote the coordinate projections
\begin{equation}
\pi_{\lambda}\colon\vec{Y}\to Y_{\lambda},
\end{equation}
where $\pi_{\lambda}(\vec{y})=y_{\lambda}$.
\end{node}

\begin{definition}
The \define{Product Topology} $\tau$ on $\vec{Y}$ is the weakest
topology that makes all coordinate projections $\pi_{\lambda}$
continuous. (We know the discrete topology is one such topology, and
we can just use Zorn's lemma.)
\end{definition}

\begin{node}
As a special case: all $(Y_{\lambda},\tau_{\lambda})$ are the same
sapce $(Y,\tau_{Y})$. The notation for the product
$\vec{Y}=Y^{\Lambda}$; so if $\Lambda=\NN$, then we have sequences of
elements of $Y$. We think of elements of $\vec{Y}$ as functions
$\vec{y}\colon\Lambda\to Y$ sending $\lambda\mapsto y_{\lambda}$.
\end{node}

\begin{node}
We have a change of notation: write $X$ instead of $\Lambda$ and
$Y^{X}$ for the set of maps $X\to Y$. We have the product topology on
the set of maps. The base for the product topology is
\begin{equation*}
\{f\mid \exists n\in\NN\ldotp f(x_{1})\in U_{1},\dots,f(x_{n})\in
U_{n}, U_{i}\in\tau_{Y}\}
\end{equation*}
So only finitely many elements $f(x_{1})$, \dots, $f(x_{n})$ belong to
open sets.
\end{node}

\begin{example}
Consider $X=\NN$ as the indexing set, and $Y=\{0,1\}$ is a 2-element
set. We give $Y$ the discrete topology $\tau=\powerset{Y}$. Then
$Y^{X}=\{0,1\}^{\NN}$ is the set of maps. This is the set of sequences
$y_{1}$, $y_{2}$, \dots, where each $y_{k}\in\{0,1\}$.

Now, consider the Cantor set formed by $K_{1}=[0,1]$,
$K_{2}=[0,1/3]\cup[2/3,1]$, \dots. Then
\begin{equation}
K=\bigcap_{j\in\NN}K_{j}
\end{equation}
is the Cantor set. We can assign a finite binary sequence (of length
$j$) to each component of $K_{j}$. This gives us a bijection
\begin{equation}
\{0,1\}^{\NN}\to(K,\tau_{K})
\end{equation}
where $\tau_{K}$ is the relative topology of the Cantor set relative
to the standard topology on $\RR$. Moreover, this is a homeomorphism!
\end{example}

\subsection{Pointwise convergence}

\begin{theorem}
Let $f_{1}\colon X\to Y$, $f_{2}\colon X\to Y$, \dots, be a sequence
of maps in $Y^{X}$. We give $Y^{X}$ the product topology. Then
$f_{n}\to f$ converges in the product topology on $Y^{X}$ if and only
if for every $x\in X$ we see the sequence $f_{n}(x)\to f(x)$ converges
in $Y$.
\end{theorem}

% p.5